\section{Das Modell}\label{sec:modell}

Dieser Abschnitt beschreibt QAM unabhängig von jeder Oberfläche und jedem Anwendungsprojekt. Die Werkzeuge der folgenden Abschnitte
(\LaTeX-Paket, Lean-Prüfer, Datenbank, Maschine) sind Umsetzungen dieses Modells; sie wenden dieselben Regeln an.

\subsection{Akten, Frames und Abbildungen}

\begin{definition}[Akte]
Eine \emph{Akte} ist eine endliche, geordnete Folge von Einträgen \(e_1,\dots,e_n\). Sie wird nur am Ende erweitert (\emph{append-only});
bestehende Einträge werden nie geändert. Eine Korrektur ist ein neuer Eintrag, der den alten ausdrücklich ersetzt.
\end{definition}

Ein \emph{Frame} ist ein Name für einen Beschreibungsrahmen: eine Konvention, ein Bezugssystem, eine Modellannahme, ein Sektor. Frames
tragen keine innere Struktur; das Modell kennt nur ihre Namen und erklärte Abbildungen zwischen ihnen. Eine \emph{Abbildung}
\(g\colon A\to B\) ist ein deklarierter Name mit Quell- und Ziel-Frame; zwei Abbildungen können zu einer dritten komponiert werden. Ein
ausgezeichneter Frame \(*\) steht für konventionsfreie Aussagen (eine reine Zahl, einen Satz der Mathematik) und ist mit jedem Frame
verträglich. Braucht eine Anwendung mehrere unabhängige Konventionsachsen, erhält jede Kombination einen eigenen Frame und es werden
Abbildungen zwischen ihnen erklärt.

\subsection{Die registrierte Aussage}

\begin{definition}[Eintrag]\label{def:eintrag}
Ein Aussageeintrag ist ein Tupel
\[
e=(\mathit{kind},\ \mathit{key},\ \mathit{frame},\ \mathit{status},\ \mathit{inputs},\ \mathit{premises},\ \mathit{statement})
\]
mit
\begin{itemize}[nosep]
\item \(\mathit{kind}\in\{\texttt{input},\texttt{claim},\texttt{thesis},\texttt{derive},\texttt{convert}\}\);
\item einem in der Akte eindeutigen Schlüssel \(\mathit{key}\);
\item einem deklarierten Frame;
\item einem Status aus der total geordneten Menge
\[
\st{open}<\st{reading}<\st{conditional}<\st{computed}<\st{proved};
\]
\item der Menge \(\mathit{inputs}\) der eingesetzten Werte, auf denen \(e\) ruht;
\item der Liste der Prämissen, das heißt früherer Einträge, aus denen \(e\) gewonnen wurde (bei \texttt{convert} zusätzlich die Abbildung);
\item optional einem Wortlaut.
\end{itemize}
\end{definition}

Die fünf Arten haben feste Bedeutungen. Eine \emph{Eingabe} ist ein von außen eingesetzter Wert: ein Messwert, ein Literaturwert, ein
angepasster Parameter. Ihre Eingabemenge enthält sie selbst. Eine \emph{Behauptung} ist eine Aussage ohne Prämissen in der Akte, etwa ein
zitierter Satz. Eine \emph{These} ist eine Behauptung, deren Wortlaut Teil ihrer Identität ist; mit ihr beginnt typischerweise eine Arbeit.
Eine \emph{Herleitung} gewinnt eine Aussage aus Prämissen im selben Frame, eine \emph{Überführung} bringt eine Aussage über eine
deklarierte Abbildung in einen anderen Frame. Daneben gibt es zwei statuslose Prüfaussagen: den \emph{Vergleich} zweier Einträge und die
\emph{Unabhängigkeitserklärung}, die behauptet, dass ein Eintrag auf keiner Eingabe ruht.

Die Statusstufen bedeuten: \st{proved} bewiesen; \st{computed} durch eine Rechnung gestützt, aber nicht bewiesen; \st{conditional} unter
benannten, nicht erfüllten Bedingungen gültig; \st{reading} eine Lesart oder Zuordnung, etwa die Identifikation eines mathematischen
Objekts mit einer physikalischen Größe; \st{open} offen. Der Status wird vom Autor vergeben. Das Modell prüft nicht, ob ein als
\st{proved} deklarierter Satz bewiesen ist, sondern nur, ob ein Schluss stärker deklariert ist als das, worauf er ruht.

\subsection{Regeln}

\begin{definition}[Wohlgeformtheit]\label{def:regeln}
Eine Akte ist \emph{wohlgeformt}, wenn für jeden Eintrag \(e_i\) gilt:
\begin{enumerate}[nosep,label=(R\arabic*)]
\item \emph{Deklaration vor Gebrauch.} Jeder Schlüssel, Frame und jede Abbildung, die \(e_i\) nennt, ist in \(e_1,\dots,e_{i-1}\) deklariert;
kein Schlüssel ist doppelt.
\item \emph{Keine Beförderung.} Ist \(e_i\) eine Herleitung oder Überführung mit Prämissen \(p_1,\dots,p_m\), so ist
\(\mathit{status}(e_i)\le\min_j\mathit{status}(p_j)\).
\item \emph{Eingaben bleiben erhalten.} \(\mathit{inputs}(e_i)=\bigcup_j\mathit{inputs}(p_j)\); eine Eingabe hat \(\mathit{inputs}=\{\mathit{key}\}\),
eine Behauptung und eine These haben \(\mathit{inputs}=\emptyset\). Eine Unabhängigkeitserklärung für \(e\) ist nur zulässig, wenn
\(\mathit{inputs}(e)=\emptyset\).
\item \emph{Keine Frame-Vermischung.} Die Prämissen einer Herleitung und die beiden Seiten eines Vergleichs liegen im selben Frame
(oder in \(*\)). Eine Überführung \(\mathrm{convert}(g,p)\) mit \(g\colon A\to B\) verlangt \(\mathit{frame}(p)=A\) und hat den Frame \(B\).
\end{enumerate}
\end{definition}

Die drei inhaltlichen Regeln entsprechen den drei Fehlerklassen der Einleitung: (R2) verhindert die Beförderung, (R3) verhindert, dass ein
eingesetzter Wert als Ergebnis gilt, (R4) verhindert die Frame-Vermischung. (R1) macht die Prüfung in einem einzigen Durchlauf
möglich.

Zwei Folgerungen sind für alles Weitere zentral. Weil (R2) und (R3) bei jedem Schritt gelten, gelten sie auch über beliebig lange Ketten:
Ein Schluss ist höchstens so stark wie der schwächste Eintrag, den man von ihm aus über Prämissen erreicht, und er trägt jede Eingabe,
die irgendwo auf diesen Wegen eingesetzt wurde. Diese Aussagen sind in Lean bewiesen (Abschnitt~\ref{sec:ebeneA}).

\subsection{Offene Bedingungen}\label{sec:offen}

Aus den gespeicherten Feldern lässt sich für jede Aussage berechnen, worauf sie letztlich ruht. Wir verwenden in diesem Artikel folgende
Sprechweise:

\begin{definition}[Offene Bedingungen einer Aussage]\label{def:offen}
Die \emph{offenen Bedingungen} einer Aussage \(e\) sind die Einträge, die von \(e\) aus über Prämissen erreichbar sind und einen Status
unterhalb von \st{proved} haben, zusammen mit \(\mathit{inputs}(e)\).
\end{definition}

Die Definition ist eine Lesart der gespeicherten Struktur, kein zusätzliches Feld; die Datenbankabfrage \texttt{premises}
(Abschnitt~\ref{sec:datenbank}) liefert genau diese erreichbaren Einträge mit ihrem Status. Nach (R2) ist eine Aussage mit einer offenen
Bedingung selbst nicht \st{proved}, und nach (R3) bleibt jede eingesetzte Größe sichtbar. Eine offene Bedingung kann deshalb nicht dadurch
verschwinden, dass weitere Schlüsse auf ihr aufbauen: Sie wird an jeden Nachfolger \emph{vererbt}. In diesem Sinn sind offene
Bedingungen Objekte des Abhängigkeitsgraphen und nicht Anmerkungen im Text. Abschnitt~\ref{sec:guards} beschreibt, wie dieselbe Idee
auf der Ebene der Ausführung umgesetzt ist.

\subsection{Identität durch Hash-Kette}

\begin{definition}[Hash eines Eintrags]
Der Hash eines Eintrags ist
\[
h(e)=\mathrm{MD5}\bigl(\mathit{kind};\mathit{key};\mathit{frame};\mathit{status};\mathit{inputs};\mathit{sources};\mathit{statement}\bigr),
\]
wobei \(\mathit{sources}\) die Liste der Hashes der Prämissen ist. Der Hash der Akte ist der Hash über die Folge aller Eintragshashes.
\end{definition}

Weil die Hashes der Prämissen in den Hash eines Schlusses eingehen, nennt jeder Schluss die exakten Fassungen dessen, worauf er ruht.
Ändert sich ein Wort einer These, ändern sich alle Hashes, die von ihr abhängen; unabhängige Einträge behalten ihre. Gleiche Hashes
bedeuten gleiche Einträge, nicht wahre. MD5 dient hier dem Erkennen versehentlicher Änderungen; für eine Veröffentlichung kommen
SHA-256-Summen über Quelle, Ausgabedokument und Akte hinzu.

\subsection{Was das Modell nicht enthält}

Das Modell kennt keine Semantik der Aussagen. Es weiß nicht, was eine Formel bedeutet, ob eine Rechnung richtig ist oder ob ein Frame
sinnvoll gewählt ist. Abbildungen haben keinen Status: Dass eine Abbildung zwischen zwei Frames physikalisch oder mathematisch
gerechtfertigt ist, muss als eigene Aussage registriert werden, wenn es eine Bedingung ist. Eine Behauptung, die ohne Prämissen als
\st{proved} deklariert wird, wird als solche akzeptiert. Ein angepasster Wert, der unter neuem Namen als Behauptung deklariert wird, wird nicht
als Eingabe erkannt. Diese Grenzen sind bewusst: Sie halten das Modell klein genug, dass seine Regeln vollständig mechanisch prüfbar sind.

\subsection{Die Objektebene}

Unterhalb der Buchführung legt QAM acht Objektarten als markierte Tupel über der gewöhnlichen Mengenlehre fest: Zustand, Schicht,
Operator, Auslesung, Zweig, Schale, Zeuge und Tensor, \(\langle0,\dots\rangle\) bis \(\langle7,\dots\rangle\). Ein eingesetzter Wert ist die
Auslesung eines Zeugen \(\langle6,u\rangle\), eine Wahl während einer Konstruktion ein Zweig; Größen tragen ihren Frame als vorderen Index,
\({}^{A}{}_{B}\mathcal T^{i}{}_{j}\). Die Objektebene fügt keine Axiome hinzu, ist in Lean formalisiert (Anhang~\ref{app:lean}) und liefert
die Bausteine der Maschinenaufträge, etwa \texttt{operator} und \texttt{state}.
